% GENERATED FILE. Edit canonical lessons, then run npm run generate:books.
% canonical-source-hash: fnv1a64:9cad86c07d9e8f1b
% canonical-lessons: ZH-C203-zu, ZH-C203-fangdong, ZH-W203-dong-east, ZH-C203-fangzu, ZH-C203-yajin, ZH-C203-hetong

\chapter{Renting a Flat}
\label{ch:zh-renting-a-flat}
\hlchaptermodality{203}

\begin{chapteropening}
\textbf{By the end of this chapter:} \emph{I can talk about renting a flat: the landlord, the rent, the deposit and the contract.}

Complete the last lesson of chapter 203: I can talk about renting a flat: the landlord, the rent, the deposit and the contract.
\end{chapteropening}

\section[zū]{zū — to rent}

\label{lesson:ZH-C203-zu}



\begin{quote}
Before the new one: say the Mandarin for customer, then the Mandarin for checkout.
\end{quote}

\subsection*{You'll want to know: zū}
\emph{zū} — \textquotedblleft{}to rent\textquotedblright{}.

\begin{quote}
\emph{wǒ xiǎng zū yí ge fángzi} — I want to rent a flat.
\end{quote}

Its characters come later; for now it is a word for the ear.

\begin{cousinweb}
\emph{zū} is inherited from Middle Chinese \emph{tsu} (Guangyun, in Baxter's transcription), glossed 'to store up' and 'a tax'. Renting is not given there, so how it came is not settled.
\end{cousinweb}

\begin{grammarlens}[title={What you've built}]
One more word for talking about plans.
\end{grammarlens}

\subsection*{Your turn}
\begin{itemize}
\raggedright
  \item \emph{Say it:} \emph{zū}
  \item \emph{Say it:} \emph{zū}, once more
  \item \emph{Say it:} \emph{wǒ xiǎng zū yí ge fángzi} — I want to rent a flat.
  \item \emph{Recall:} say the Mandarin for customer, then the Mandarin for checkout, then say \emph{zū} again
\end{itemize}

\begin{tcolorbox}[breakable,colback=teal!4,colframe=teal!35!black,title={Before you move on}]
What does zū mean? (\textquotedblleft{}To rent\textquotedblright{}.) Say it once more.
\end{tcolorbox}
\section[\texorpdfstring{\zh{房东} (fángdōng)}{房东 (fángdōng)}]{\zh{房东} (fángdōng) — landlord; landlady}

\label{lesson:ZH-C203-fangdong}



\begin{quote}
Before the new one: say the Mandarin for checkout, then the Mandarin for to rent.
\end{quote}

\subsection*{You'll want to know: \zh{房东}}
\textbf{\zh{房东}} — \emph{fángdōng} — \textquotedblleft{}landlord; landlady\textquotedblright{}.

\begin{quote}
\textbf{\zh{我的房东人很好}} \emph{wǒ de fángdōng rén hěn hǎo} — My landlord is a nice person.
\end{quote}

The new character, in pieces: \zh{东} — five strokes: a bar, a stroke that falls to the left and turns right, a vertical that hooks to the left, then a short stroke falling to the left and a dot. The next lesson puts it in your hand.

\begin{cousinweb}
\textbf{\zh{房东}} (\emph{fángdōng}) is built from \zh{房} 'house, room' and \zh{东} 'east', both inherited from Middle Chinese, \emph{bjang} and \emph{tuwng} (Guangyun, in Baxter's transcription). How the pair came to mean the owner who lets a room is not settled.
\end{cousinweb}

\begin{grammarlens}[title={What you've built}]
One more word for talking about plans.
\end{grammarlens}

\subsection*{Your turn}
\begin{itemize}
\raggedright
  \item \emph{Say it:} \emph{fángdōng}
  \item \emph{Say it:} \emph{fángdōng}, once more
  \item \emph{Say it:} \emph{wǒ de fángdōng rén hěn hǎo} — My landlord is a nice person.
  \item \emph{Recall:} say the Mandarin for checkout, then the Mandarin for to rent, then say \emph{fángdōng} again
\end{itemize}

\begin{tcolorbox}[breakable,colback=teal!4,colframe=teal!35!black,title={Before you move on}]
What does \zh{房东} mean? (\textquotedblleft{}Landlord; landlady\textquotedblright{}.) Say it once more.
\end{tcolorbox}
\section[\texorpdfstring{\zh{东} (dōng)}{东 (dōng)}]{\zh{东} — a bar, a turn, a hook and two short strokes}

\label{lesson:ZH-W203-dong-east}



\begin{quote}
The word first, then the character inside it.

\begin{itemize}
\raggedright
  \item \emph{Recall:} say \emph{fángdōng} — \textbf{R1}, one lesson back, before it decays
\end{itemize}

\begin{quote}
\zh{房} \zh{车}
\end{quote}

\emph{Read it:} each one aloud, with nothing said first — \emph{fáng}, \emph{chē}
\end{quote}

\subsection*{Script — a bar, a turn, a hook and two short strokes}
\hlblockfigure{\detokenize{figures/ZH-W203-dong-east-filmstrip.pdf}}{How \zh{东} is written, stroke by stroke}

\begin{quote}
\zh{东}
\end{quote}

\textbf{One piece, five strokes, top to bottom.} Count the strokes in their turns.

\textbf{First, a bar} across the upper middle.

\textbf{Second, a stroke from above the bar,} falling down to the left and turning to run right.

\textbf{Third, a vertical} down the middle, from above the turned stroke and down through it, hooking to the left at its foot.

\textbf{Last, two short strokes at the foot:} one falling to the left, then a dot on the right.

1 + 1 + 1 + 2 = five.

\subsection*{Your turn}
\begin{itemize}
\raggedright
  \item \emph{Write it:} \zh{东}, all five, with the model in front of you
  \item \emph{Cover:} write it again from memory, and count the pieces, not the strokes
\end{itemize}

\begin{tcolorbox}[breakable,colback=teal!4,colframe=teal!35!black,title={Before you move on}]
How many strokes in \zh{东}? (\textbf{Five}.) Which stroke hooks? (\textbf{The vertical down the middle}.) What comes last? (\textbf{The dot on the right}.)

Source: \href{https://raw.githubusercontent.com/chanind/hanzi-writer-data/68d10a4b21150cae5e1ebbd223eed289cf32d90c/data/\%E4\%B8\%9C.json}{Hanzi Writer Data \zh{东}, pinned stroke paths and medians}.
\end{tcolorbox}
\section[fángzū]{fángzū — rent for a home}

\label{lesson:ZH-C203-fangzu}



\begin{quote}
Before the new one: say the Mandarin for to rent, then the Mandarin for landlord.
\end{quote}

\subsection*{You'll want to know: fángzū}
\emph{fángzū} — \textquotedblleft{}rent for a home\textquotedblright{}.

\begin{quote}
\emph{zhèr de fángzū tài guì le} — The rent here is too high.
\end{quote}

Its characters come later; for now it is a word for the ear.

\begin{cousinweb}
\emph{fángzū} is built from \zh{房} 'house' and \emph{zū} 'rent', both inherited from Middle Chinese, \emph{bjang} and \emph{tsu} (Guangyun, in Baxter's transcription), which glosses the second 'a tax'.
\end{cousinweb}

\begin{grammarlens}[title={What you've built}]
One more word for talking about plans.
\end{grammarlens}

\subsection*{Your turn}
\begin{itemize}
\raggedright
  \item \emph{Say it:} \emph{fángzū}
  \item \emph{Say it:} \emph{fángzū}, once more
  \item \emph{Say it:} \emph{zhèr de fángzū tài guì le} — The rent here is too high.
  \item \emph{Recall:} say the Mandarin for to rent, then the Mandarin for landlord, then say \emph{fángzū} again
\end{itemize}

\begin{tcolorbox}[breakable,colback=teal!4,colframe=teal!35!black,title={Before you move on}]
What does fángzū mean? (\textquotedblleft{}Rent for a home\textquotedblright{}.) Say it once more.
\end{tcolorbox}
\section[yājīn]{yājīn — a deposit paid in advance}

\label{lesson:ZH-C203-yajin}



\begin{quote}
Before the new one: say the Mandarin for landlord, then the Mandarin for rent for a home.
\end{quote}

\subsection*{You'll want to know: yājīn}
\emph{yājīn} — \textquotedblleft{}a deposit paid in advance\textquotedblright{}.

\begin{quote}
\emph{yājīn shì liǎng ge yuè de fángzū} — The deposit is two months' rent.
\end{quote}

Its characters come later; for now it is a word for the ear.

\begin{cousinweb}
\emph{yājīn} is built from \emph{yā} 'a pledge' and \emph{jīn} 'gold, money', both inherited from Middle Chinese, \emph{'aep} and \emph{kim} (Guangyun, in Baxter's transcription). The dictionary glosses the first as putting one's mark on a document, so how it came to mean a pledge is not settled.
\end{cousinweb}

\begin{grammarlens}[title={What you've built}]
One more word for talking about plans.
\end{grammarlens}

\subsection*{Your turn}
\begin{itemize}
\raggedright
  \item \emph{Say it:} \emph{yājīn}
  \item \emph{Say it:} \emph{yājīn}, once more
  \item \emph{Say it:} \emph{yājīn shì liǎng ge yuè de fángzū} — The deposit is two months' rent.
  \item \emph{Recall:} say the Mandarin for landlord, then the Mandarin for rent for a home, then say \emph{yājīn} again
\end{itemize}

\begin{tcolorbox}[breakable,colback=teal!4,colframe=teal!35!black,title={Before you move on}]
What does yājīn mean? (\textquotedblleft{}A deposit paid in advance\textquotedblright{}.) Say it once more.
\end{tcolorbox}
\section[hétong]{hétong — contract; written agreement}

\label{lesson:ZH-C203-hetong}



\begin{quote}
Before the new one: say the Mandarin for rent for a home, then the Mandarin for a deposit paid in advance.
\end{quote}

\subsection*{You'll want to know: hétong}
\emph{hétong} — \textquotedblleft{}contract; written agreement\textquotedblright{}.

\begin{quote}
\emph{fángdōng de hétong nǐ kànle ma} — Have you read the landlord's contract?
\end{quote}

Its characters come later; for now it is a word for the ear.

\begin{cousinweb}
\emph{hétong} is built from \emph{hé} 'to join' and \zh{同} 'together', both inherited from Middle Chinese, \emph{hop} and \emph{duwng} (Guangyun, in Baxter's transcription), which glosses the second 'alike' and 'together'. The dictionary already pairs the two to gloss the first, in the sense of joining; a written agreement is not given there, so how the word came to it is not settled. In the word the second loses its tone.
\end{cousinweb}

\begin{grammarlens}[title={What you've built}]
One more word for talking about plans.
\end{grammarlens}

\subsection*{Your turn}
\begin{itemize}
\raggedright
  \item \emph{Say it:} \emph{hétong}
  \item \emph{Say it:} \emph{hétong}, once more
  \item \emph{Say it:} \emph{fángdōng de hétong nǐ kànle ma} — Have you read the landlord's contract?
  \item \emph{Recall:} say the Mandarin for rent for a home, then the Mandarin for a deposit paid in advance, then say \emph{hétong} again
\end{itemize}

\begin{tcolorbox}[breakable,colback=teal!4,colframe=teal!35!black,title={Before you move on}]
What does hétong mean? (\textquotedblleft{}Contract; written agreement\textquotedblright{}.) Say it once more.
\end{tcolorbox}
